% =====================================================================
% Drop-in replacement for the Simulation Results and Discussion section
% of OCP.tex. Figure paths assume figures live under Figures/GA/Results/.
% Adjust the \includegraphics paths if the project tree differs.
% =====================================================================

\section{Simulation Results and Discussion}
\label{sec: sim_results}
The simulation campaign comprises $80$ independent runs of the guided GA at each design cell, where a cell is one combination of camera count $k\in\{6,7,8\}$, target type $\mathrm{TT}\in\{\text{UAV},\text{UGV}\}$, grid mode $\mathrm{GM}\in\{\text{Uniform},\text{Normal}\}$, and warm-start state. Across all reported cells the GA executes the same evaluation budget of $N_{\mathrm{pop}}+G_{\max}\!\cdot\!N_C = 100 + 100\!\times\!150 = 15{,}100$ cost evaluations, so reported differences in solution quality are attributable to the search dynamics rather than to differing computational budgets. The representative seven-camera configuration is the headline case because it matches the OptiTrack rig installed in the MechatronicSystems.Group laboratory and is therefore the configuration that the physical verification framework of Section~\ref{sec:verification} will eventually scrutinise. Six- and eight-camera results are reported as ablations only, to demonstrate how cost scales with the camera budget. All quoted costs are the combined normalised objective $J$ defined in Equation~\eqref{eq:combined-cost} with $w_{\mathrm{res}}=w_{\mathrm{occ}}=0.5$.

The remainder of this section reports four families of results. Section~\ref{sec:sim-convergence} characterises the convergence behaviour of the GA, including the effect of warm starting and elitism activation on the speed and quality of convergence. Section~\ref{sec:sim-costeval} compares the GA-optimised configurations against the OptiTrack ad-hoc baseline on the combined cost and on the constituent resolution-uncertainty and dynamic-occlusion components. Section~\ref{sec:sim-sensitivity} reports the sensitivity of the cost function to two design choices: the target-grid spacing $\Delta$ and per-camera orientation perturbations, the former defending the choice of $\Delta=1$~m and the latter demonstrating that each returned camera lies at or very near a local minimum of the cost surface. Section~\ref{sec:sim-optconfig} interprets the spatial structure of the optimal configurations and shows how the choice of target-grid modality (UAV full-volume vs UGV ground-slab) produces materially different placements.

\subsection{Convergence Behaviour}
\label{sec:sim-convergence}

Convergence curves for the seven-camera, $1$~m grid spacing case are shown in Figure~\ref{fig:convergence-7c}. Each panel splits the $40$ independent runs in that cell by start strategy: cold-start (left) and warm-start (right). Cold-start chromosomes are seeded by the guided perimeter-and-aim procedure of Algorithm~\ref{alg:init-pop} but otherwise initialised randomly; warm-start runs additionally inject a small number of elite chromosomes from previous cold-start runs into the initial population.

\begin{figure}[htbp]
    \centering
    \begin{subfigure}[t]{0.49\textwidth}
        \centering
        \includegraphics[width=\textwidth]{Figures/GA/Results/convergence_CF3_Combined_UAV_GM1_Uniform_sp100cm_7C.pdf}
        \caption{UAV target type, uniform grid.}
        \label{fig:convergence-uav-uniform}
    \end{subfigure}
    \hfill
    \begin{subfigure}[t]{0.49\textwidth}
        \centering
        \includegraphics[width=\textwidth]{Figures/GA/Results/convergence_CF3_Combined_UAV_GM2_Normal_sp100cm_7C.pdf}
        \caption{UAV target type, normal-concentrated grid.}
        \label{fig:convergence-uav-normal}
    \end{subfigure}

    \vspace{0.4em}

    \begin{subfigure}[t]{0.49\textwidth}
        \centering
        \includegraphics[width=\textwidth]{Figures/GA/Results/convergence_CF3_Combined_UGV_GM1_Uniform_sp100cm_7C.pdf}
        \caption{UGV target type, uniform grid.}
        \label{fig:convergence-ugv-uniform}
    \end{subfigure}
    \hfill
    \begin{subfigure}[t]{0.49\textwidth}
        \centering
        \includegraphics[width=\textwidth]{Figures/GA/Results/convergence_CF3_Combined_UGV_GM2_Normal_sp100cm_7C.pdf}
        \caption{UGV target type, normal-concentrated grid.}
        \label{fig:convergence-ugv-normal}
    \end{subfigure}

    \caption[Convergence of the combined cost across $40$ runs per cell.]{Convergence trajectories of the seven-camera GA under each (target~type, grid~mode) cell at $\Delta=1$~m. Each panel separates cold-start (left) and warm-start (right) runs: the solid green band traces the median best-cost with the $Q25$--$Q75$ inter-quartile ribbon, the dashed blue curve the median top-10 population average, the deep-red trace highlights the run achieving the lowest final cost, and the dash-dot line marks the per-panel best.}
    \label{fig:convergence-7c}
\end{figure}

Three observations are common to all four cells. First, the median best-cost curve drops sharply during the first $\sim 20$ generations and then enters a long shallow refinement phase that continues until $G_{\max}$. For the UAV/Uniform cold-start panel of Figure~\ref{fig:convergence-uav-uniform}, the median falls from an initial $0.115$ to $0.047$ over the first $50$ generations, a $59\%$ reduction, and continues to drift down to $0.040$ by generation $150$, a further $15\%$. This is consistent with the design of the operator schedule: during the exploratory phase ($g<i t_e = 50$) the offspring fully replace the parents and the median can both rise and fall; once elitism activates the merged-and-truncated population guarantees that the best individual cannot regress, producing the characteristic flattening of the median trajectory. The same pattern is observed in the UGV/Uniform case (Figure~\ref{fig:convergence-ugv-uniform}), where the cold-start median falls from $0.086$ at initialisation to $0.027$ by generation $150$, and in both normal-concentrated cells.

Second, the top-10 average (median across runs) tracks above the best-cost curve throughout the run but converges towards it. In the UAV/Uniform cold-start panel the top-10 curve sits roughly $0.010$ above the best curve for the first hundred generations and closes the gap to $\sim 0.005$ by generation $150$. This is the population-level evidence of convergence that elitism alone cannot certify: the elite individual could in principle remain unique with the rest of the population unchanged, but the observed downward drift of the top-10 average shows that the genetic operators are actively raising the floor of the population, not merely preserving a lucky outlier. This is the same evidence that Aissaoui et al.\ used to defend the guided GA over an unseeded GA in \parencite{aissaoui_designing_2018}: the population average, not just the best, must improve over generations.

Third, the population-diversity decay shown in Figure~\ref{fig:diversity-7c} confirms that the GA enters its exploitation phase well before elitism activates. Diversity, computed as the median per-gene standard deviation across the population, collapses from $\approx 2.0$~(a.u.) at $g=0$ to $\approx 0.15$ by $g=50$, after which it remains essentially flat for the remaining $100$ generations. The best-cost median, by contrast, continues to drop over the same interval. This decoupling is the signature of a healthy elitist GA: when diversity has saturated, mutation and the camera-repair operator (Section~\ref{sec:mutation}) are still finding improving moves on the boundary of the converged cluster. If diversity had collapsed before the best-cost curve flattened, the algorithm would be exhibiting premature convergence; if diversity had remained high without the best cost improving, the algorithm would be exploring without exploiting. Neither pathology is observed here.

\begin{figure}[htbp]
    \centering
    \begin{subfigure}[t]{0.49\textwidth}
        \centering
        \includegraphics[width=\textwidth]{Figures/GA/Results/diversity_CF3_Combined_UAV_GM1_Uniform_sp100cm_7C.pdf}
        \caption{UAV, uniform grid.}
        \label{fig:diversity-uav}
    \end{subfigure}
    \hfill
    \begin{subfigure}[t]{0.49\textwidth}
        \centering
        \includegraphics[width=\textwidth]{Figures/GA/Results/diversity_CF3_Combined_UGV_GM1_Uniform_sp100cm_7C.pdf}
        \caption{UGV, uniform grid.}
        \label{fig:diversity-ugv}
    \end{subfigure}
    \caption[Population-diversity decay over generations.]{Median per-gene population diversity (solid red) and median normalised best cost (dashed green) across all $40$ runs in the (target~type, grid~mode) cell. Diversity collapses by $g\!\approx\!50$ whilst the best cost continues to improve, evidencing healthy exploitation under elitism without premature convergence.}
    \label{fig:diversity-7c}
\end{figure}

The effect of warm starting is consistent across all four cells. Table~\ref{tab:warm-vs-cold} summarises the per-cell statistics over the $20$ runs of each start strategy at $k=7$ and $\Delta=1$~m. Warm starting both reduces the median final cost and tightens its dispersion. In the UAV/Uniform cell the cold-start median final cost is $0.0478$ with a standard deviation of $0.0052$, whereas the warm-start median is $0.0375$ with a standard deviation of $0.0019$, a $22\%$ reduction in median and a $63\%$ reduction in run-to-run variability. The same direction of effect holds in all other cells, with the warm-start median cost lying $19$--$30\%$ below the cold-start median and the warm-start standard deviation $44$--$72\%$ below the cold-start standard deviation. The interpretation is that warm starting effectively bypasses the exploratory descent of Figure~\ref{fig:convergence-7c}: the initial population already sits near a known low-cost basin, and the remaining $150$ generations are spent in localised refinement rather than in the early-generation phase of escaping the random initial scatter. This mirrors the rationale Aissaoui et al.\ gave for their guided initialisation in \parencite{aissaoui_designing_2018}, but takes it one step further by using the GA's own elite history as the seed, rather than a domain-derived trajectory.

\begin{table}[htbp]
\centering
\caption[Best-cost statistics for warm- vs cold-start runs, $k=7$, $\Delta=1$~m.]{Best-cost statistics for warm-start versus cold-start initialisation at the representative seven-camera, $\Delta=1$~m configuration. Each cell summarises $n=20$ independent runs.}
\label{tab:warm-vs-cold}
\small
\begin{tabular}{l l l l l l l}
\toprule
TT & GM & Start & Min & Median & Mean & Std \\
\midrule
UAV & Uniform & Cold & 0.0374 & 0.0478 & 0.0470 & 0.0052 \\
UAV & Uniform & Warm & 0.0330 & 0.0375 & 0.0379 & 0.0019 \\
UAV & Normal  & Cold & 0.0301 & 0.0376 & 0.0384 & 0.0046 \\
UAV & Normal  & Warm & 0.0281 & 0.0301 & 0.0305 & 0.0016 \\
UGV & Uniform & Cold & 0.0227 & 0.0270 & 0.0287 & 0.0066 \\
UGV & Uniform & Warm & 0.0179 & 0.0227 & 0.0216 & 0.0018 \\
UGV & Normal  & Cold & 0.0151 & 0.0214 & 0.0208 & 0.0027 \\
UGV & Normal  & Warm & 0.0126 & 0.0151 & 0.0157 & 0.0014 \\
\bottomrule
\end{tabular}
\end{table}

The robustness annotation embedded in each convergence panel (``$X/20$ runs within $5\%$ of panel best'') quantifies a related but distinct property: how often the algorithm lands on its own best basin, rather than on a sub-optimal local minimum. In the cold-start panels this value ranges from $1/20$ (UGV, normal, $5\%$) to $4/20$ (UGV, uniform, $20\%$); in the warm-start panels it ranges from $1/20$ to $3/20$. The fact that warm starting reduces variability (Table~\ref{tab:warm-vs-cold}) without proportionally increasing the within-$5\%$ count indicates that warm starting concentrates runs in a band that lies just outside the $5\%$ tolerance of the single best run rather than collapsing the entire run set onto the same global minimum. The cost landscape, in other words, contains a small number of competitive basins of comparable quality; each independent run reliably finds \emph{one} of them, but no single run is guaranteed to find the deepest. The implication for the practitioner is that a small number of independent restarts is required to extract the best configuration with high confidence; the cost of this is modest given the per-run wall-clock time of $\sim 20$ minutes (UAV) or $\sim 13$ minutes (UGV).

\subsection{Cost Function Evaluation}
\label{sec:sim-costeval}

This subsection compares the GA-optimised configurations against the OptiTrack ad-hoc baseline on the combined objective $J$ and on its two components, and reports how the cost scales with camera count and with the choice of target grid.

Figure~\ref{fig:cost-boxplots} presents the best-cost distribution across all runs as box-and-whisker plots, with the OptiTrack ad-hoc cost (Figure~\ref{fig:cost-uniform}) overlaid as a single marker. Two observations dominate. Firstly, the OptiTrack cost lies an order of magnitude above the GA-optimised cost across every camera count: at $k=7$ and uniform grid the OptiTrack cost is $0.2691$ against a GA median of $0.0400$, a $6.7\times$ improvement; at $k=7$ and normal grid the OptiTrack cost is $0.2393$ against a GA median of $0.0324$, a $7.4\times$ improvement. For the UGV target the corresponding ratios are $5.8\times$ (uniform, $0.1449$ vs $0.0232$) and $7.1\times$ (normal, $0.1218$ vs $0.0172$). These magnitudes are consistent with the experimental gain Rahimian and Kearney reported on their simulated-annealing-optimised configurations, where the optimised placement achieved full marker reconstruction on $86\%$ of frames against $30\%$ for the ad-hoc layout, a $\sim 2.9\times$ ratio on a different but related metric \parencite{rahimian_optimal_2017}; the larger ratio here reflects both the budget-constrained sensor model (narrower FOV, shorter range) and the composite uncertainty-plus-occlusion objective, neither of which the ad-hoc OptiTrack arrangement was designed to satisfy.

\begin{figure}[htbp]
    \centering
    \begin{subfigure}[t]{0.49\textwidth}
        \centering
        \includegraphics[width=\textwidth]{Figures/GA/Results/cost_byTT_CF3_Combined_GM1_Uniform_sp100cm_CF3.pdf}
        \caption{Uniform grid.}
        \label{fig:cost-uniform}
    \end{subfigure}
    \hfill
    \begin{subfigure}[t]{0.49\textwidth}
        \centering
        \includegraphics[width=\textwidth]{Figures/GA/Results/cost_byTT_CF3_Combined_GM2_Normal_sp100cm_CF3.pdf}
        \caption{Normal-concentrated grid.}
        \label{fig:cost-normal}
    \end{subfigure}
    \caption[Combined-cost distributions across camera counts.]{Best-cost distribution of the GA at $k\in\{6,7,8\}$ for both target types under the two grid modes, with the OptiTrack ad-hoc cost overlaid as a single marker. Boxes show the inter-quartile range, whiskers the $1.5\times$IQR extents, and the central line the median; $n$ annotations record the number of runs per box.}
    \label{fig:cost-boxplots}
\end{figure}

Secondly, the GA cost decreases monotonically with camera count, as expected from the physics. From $k=6$ to $k=8$ the UAV uniform-grid median falls from $0.045$ to $0.035$ ($-22\%$); the UGV uniform-grid median falls from $0.024$ to $0.013$ ($-46\%$). The marginal benefit of the eighth camera is smaller than that of the seventh, suggesting that the workspace is approaching coverage saturation at $k=7$ for the UGV case in particular, where every floor target can already be observed by an average of $4.65$ cameras (Table~\ref{tab:coverage-stats}). The trend is consistent with the diminishing-returns observation that Zhou et al.\ reported when sweeping camera count in their hybrid PSO+GA formulation \parencite{zhou_sensor_2024}: once each target point is covered by a triangulable pair the marginal cameras can only reduce uncertainty, not unlock new coverage.

\begin{table}[htbp]
\centering
\caption[Per-target-point coverage statistics for the GA-optimised configurations.]{Coverage statistics averaged across all $40$ runs per cell at $k=7$ and $\Delta=1$~m. The $0$-camera column reports the percentage of target points with no observing camera; the $1$-camera column reports the percentage observed by exactly one camera (which cannot be triangulated); the $2{+}$-camera column reports the percentage observed by at least two cameras (triangulable, subject to the convergence-angle bound).}
\label{tab:coverage-stats}
\small
\begin{tabular}{l l r r r r}
\toprule
TT & GM & \multicolumn{1}{c}{$0$-cam (\%)} & \multicolumn{1}{c}{$1$-cam (\%)} & \multicolumn{1}{c}{$2{+}$-cam (\%)} & Avg.\ coverage \\
\midrule
UAV & Uniform & 0.01 & 19.33 & 80.67 & 3.85 \\
UAV & Normal  & 0.01 & 12.83 & 87.17 & 4.43 \\
UGV & Uniform & 0.00 & 20.41 & 79.59 & 4.33 \\
UGV & Normal  & 0.00 & 15.80 & 84.20 & 4.65 \\
\bottomrule
\end{tabular}
\end{table}

The decomposition of the combined cost into its uncertainty and occlusion components is shown spatially in Figure~\ref{fig:costfield}. For the UAV target the GA-best layout (left panels) produces a per-point uncertainty field that is uniform across the volume with a mean of $\sim 6$~px and no high-cost lobes, whereas the OptiTrack ad-hoc layout (right panels) exhibits pronounced high-uncertainty regions at the corners of the volume where targets fall outside the redundant-coverage envelope. The corresponding per-point occlusion field tells the same story: under the GA layout the occlusion metric is uniformly below $200^\circ$ at every target point, whereas the ad-hoc layout exhibits points exceeding $400^\circ$ in the upper corners, indicating that for those targets fewer than two cameras have a triangulable view across more than half of the possible occluder orientations.

\begin{figure}[htbp]
    \centering
    \begin{subfigure}[t]{0.49\textwidth}
        \centering
        \includegraphics[width=\textwidth]{Figures/GA/Results/CostField_GAvsOpti_UAV.pdf}
        \caption{UAV full-volume.}
        \label{fig:costfield-uav}
    \end{subfigure}
    \hfill
    \begin{subfigure}[t]{0.49\textwidth}
        \centering
        \includegraphics[width=\textwidth]{Figures/GA/Results/CostField_GAvsOpti_UGV.pdf}
        \caption{UGV ground slab.}
        \label{fig:costfield-ugv}
    \end{subfigure}
    \caption[Per-point cost decomposition: GA-best vs OptiTrack ad-hoc.]{Spatial distribution of the resolution-uncertainty (top row of each panel) and dynamic-occlusion (bottom row) components of $J$ across the target grid, comparing the GA-best $k=7$ configuration with the OptiTrack ad-hoc baseline. The GA layout produces uniformly low cost across the volume; the ad-hoc layout exhibits localised high-cost regions where the placement fails to provide a triangulable view.}
    \label{fig:costfield}
\end{figure}

A complementary view is provided by the pairwise camera-baseline-angle distribution in Figure~\ref{fig:baseline-angles}. The triangulability constraint of Section~\ref{paragraph: rahimian_refinement} accepts a camera pair as a useful observer of a given target only when the convergence angle lies in $[40^\circ, 140^\circ]$. The GA-best UAV configuration places $73.2\%$ of its camera pairs in this band, against $64.8\%$ for the OptiTrack rig; under the UGV target the figures are $79.9\%$ versus $73.7\%$ respectively. The OptiTrack histogram exhibits a sharp peak at $\sim 20^\circ$, far below the lower band edge, evidencing the clustered placement visible in Figure~\ref{fig:configfov-uav}; the GA histogram shifts mass towards the $60^\circ$--$120^\circ$ region where triangulation is best conditioned. This is the same diagnostic Olague and Mohr used in \parencite{olague_optimal_2002} to defend their covariance-minimising configurations, but expressed in the angular rather than the covariance frame.

\begin{figure}[htbp]
    \centering
    \begin{subfigure}[t]{0.49\textwidth}
        \centering
        \includegraphics[width=\textwidth]{Figures/GA/Results/BaselineAngles_GAvsOpti_UAV.pdf}
        \caption{UAV target.}
        \label{fig:baseline-uav}
    \end{subfigure}
    \hfill
    \begin{subfigure}[t]{0.49\textwidth}
        \centering
        \includegraphics[width=\textwidth]{Figures/GA/Results/BaselineAngles_GAvsOpti_UGV.pdf}
        \caption{UGV target.}
        \label{fig:baseline-ugv}
    \end{subfigure}
    \caption[Pairwise baseline-angle distribution.]{Distribution of pairwise camera baseline angles for the GA-best configuration (green) and the OptiTrack ad-hoc baseline (red). Dashed vertical lines mark the triangulability bounds $[40^\circ, 140^\circ]$. The GA layout shifts pairs into the well-conditioned interior of this band; the ad-hoc layout retains a large near-parallel peak below the lower bound.}
    \label{fig:baseline-angles}
\end{figure}

The computational cost of these optimisations is reported in Figure~\ref{fig:comptime}. Per-run wall-clock time is dominated by the number of target points $|\mathcal{T}|$, not by the chromosome length: at $\Delta=1$~m the UAV grid contains $\sim 405$ points and the UGV slab $\sim 243$ points, with the UAV runs taking a median $\sim 20.8$~minutes against a UGV median of $\sim 13.6$~minutes ($1.5\times$, in close proportion to the $1.7\times$ point ratio). The variation with camera count is comparatively modest: increasing $k$ from $6$ to $8$ raises the median wall-clock by less than $30\%$, since the per-point cost evaluation is dominated by the pairwise visibility, range, and triangulability checks rather than by the $\mathcal{O}(k)$ FOV tests. The outliers above one hour visible in Figure~\ref{fig:comptime-tt-uniform} arise from a small number of pilot runs that were executed on a shared machine; they are retained in the figure for completeness rather than excluded.

\begin{figure}[htbp]
    \centering
    \begin{subfigure}[t]{0.49\textwidth}
        \centering
        \includegraphics[width=\textwidth]{Figures/GA/Results/comptime_byTT_GM1_Uniform_sp100cm.pdf}
        \caption{Split by target type (uniform grid).}
        \label{fig:comptime-tt-uniform}
    \end{subfigure}
    \hfill
    \begin{subfigure}[t]{0.49\textwidth}
        \centering
        \includegraphics[width=\textwidth]{Figures/GA/Results/comptime_byCF_GM1_Uniform_sp100cm.pdf}
        \caption{Pooled across target types (uniform grid).}
        \label{fig:comptime-cf-uniform}
    \end{subfigure}
    \caption[Per-run wall-clock time.]{Distribution of per-run computation time across the $40$ runs of each cell at $\Delta=1$~m on $24$ parallel workers. The UAV runs take roughly $1.5\times$ longer than the UGV runs due to the larger target grid; the dependence on camera count is mild.}
    \label{fig:comptime}
\end{figure}

\subsection{Parameter Sensitivity Analysis}
\label{sec:sim-sensitivity}

This subsection reports two distinct sensitivities. The first is the sensitivity of the cost function itself to the target-grid spacing $\Delta$, which determines how faithfully the discrete cost approximates the continuous integral over the workspace. The second is the sensitivity of the optimised cost to small perturbations of each camera's orientation, which establishes whether the GA's returned configurations sit at true local minima of the cost surface or merely at saddle-like configurations the operator schedule happened to abandon.

\paragraph{Target-grid spacing:}
Figure~\ref{fig:spacing-sensitivity-uav} sweeps the grid spacing from $\Delta=0.1$~m (the finest grid practically tractable for the UAV volume, $\approx 200{,}000$ target points) to $\Delta=3$~m. The GA-best chromosome and the OptiTrack ad-hoc chromosome are held fixed; only the evaluation grid changes. For the GA-best configuration the resolution-uncertainty raw mean is stable to within $\pm 2\%$ across the entire range, indicating that the back-projected pyramid intersection at well-separated target points is a low-curvature function of position and is therefore well approximated even by coarse grids. The dynamic-occlusion raw mean is more sensitive: at $\Delta\geq 1$~m the metric begins to drift upward by $\sim 50\%$, because the coarser grid begins to alias the angular-sector partition. The normalised combined CF3 inherits both behaviours but in milder form: at $\Delta=1$~m the combined cost deviates from the $\Delta=0.1$~m reference by $\sim 4\%$ for the GA-best chromosome and $\sim 5\%$ for the OptiTrack chromosome.

\begin{figure}[htbp]
    \centering
    \begin{subfigure}[t]{\textwidth}
        \centering
        \includegraphics[width=\textwidth]{Figures/GA/Results/Sensitivity/UAV/cost_vs_spacing.png}
        \caption{Raw component and normalised combined cost versus grid spacing $\Delta$ for the UAV full-volume.}
        \label{fig:spacing-uav-cost}
    \end{subfigure}

    \vspace{0.4em}

    \begin{subfigure}[t]{\textwidth}
        \centering
        \includegraphics[width=\textwidth]{Figures/GA/Results/Sensitivity/UAV/relative_deviation.png}
        \caption{Relative deviation of each metric from the $\Delta=0.1$~m reference.}
        \label{fig:spacing-uav-dev}
    \end{subfigure}

    \vspace{0.4em}

    \begin{subfigure}[t]{\textwidth}
        \centering
        \includegraphics[width=\textwidth]{Figures/GA/Results/Sensitivity/UAV/eval_time.png}
        \caption{Per-evaluation wall-clock time and target-point count versus $\Delta$.}
        \label{fig:spacing-uav-time}
    \end{subfigure}

    \caption[Grid-spacing sensitivity, UAV target type.]{Sensitivity of the cost components and per-evaluation runtime to the target-grid spacing $\Delta$. The combined CF3 deviation from the finest-grid reference remains below $5\%$ for $\Delta\leq 1$~m, whilst the per-evaluation cost drops by more than two orders of magnitude across the same range, defending the choice of $\Delta=1$~m for the GA's inner cost evaluation.}
    \label{fig:spacing-sensitivity-uav}
\end{figure}

The corresponding evaluation-time plot of Figure~\ref{fig:spacing-uav-time} makes the trade-off explicit. A single cost evaluation at $\Delta=0.1$~m takes $\approx 38$~seconds, whilst the same evaluation at $\Delta=1$~m takes $\approx 0.15$~seconds, a $\sim 250\times$ speedup. Multiplied across the $15{,}100$ evaluations per run, the $\Delta=0.1$~m setting would extend each run from the observed $\sim 20$~minutes to over $150$~hours, which is infeasible at the campaign scale of this study. The choice of $\Delta=1$~m therefore trades a $\sim 4\%$ approximation error against a tractable run budget, with the residual error absorbed by the same monotonic mapping from cost to GA-best as the finer grids. The corresponding UGV sweep (Figure~\ref{fig:spacing-sensitivity-ugv}, sampled at $\Delta=0.25$--$1$~m as the slab geometry tolerates the coarser bound) confirms the same picture: combined CF3 stable to within $\pm 3\%$ across the entire sweep, with the dominant cost being target-point count.

\begin{figure}[htbp]
    \centering
    \includegraphics[width=\textwidth]{Figures/GA/Results/Sensitivity/UGV/cost_vs_spacing.png}
    \caption[Grid-spacing sensitivity, UGV target type.]{Raw and normalised cost versus grid spacing on the UGV ground slab. The combined cost is essentially flat across the swept range, reflecting the smaller and lower-dimensional target set.}
    \label{fig:spacing-sensitivity-ugv}
\end{figure}

\paragraph{Per-camera orientation sensitivity:}
A different question is whether each camera in the GA-best configuration sits at a local minimum of the cost surface with respect to its own orientation. Figure~\ref{fig:angle-sensitivity} probes this by perturbing camera $1$'s roll, pitch, and yaw individually across $\pm 45^\circ$ while holding the remaining six cameras fixed, and re-evaluating the combined cost at each perturbation.

\begin{figure}[htbp]
    \centering
    \begin{subfigure}[t]{0.49\textwidth}
        \centering
        \includegraphics[width=\textwidth]{Figures/GA/Results/AngleSensitivity_UAV_cam1.pdf}
        \caption{UAV target type.}
        \label{fig:angle-uav}
    \end{subfigure}
    \hfill
    \begin{subfigure}[t]{0.49\textwidth}
        \centering
        \includegraphics[width=\textwidth]{Figures/GA/Results/AngleSensitivity_UGV_cam1.pdf}
        \caption{UGV target type.}
        \label{fig:angle-ugv}
    \end{subfigure}
    \caption[Per-camera orientation sensitivity at the GA-best.]{Sensitivity of the combined cost to one-at-a-time roll, pitch, and yaw perturbations of camera $1$ in the GA-best $k=7$ configuration. Each curve traces $J$ as the named Euler angle is perturbed by $\pm 45^\circ$ from the optimised value (vertical dashed line); the marker at $0^\circ$ indicates the GA's returned cost.}
    \label{fig:angle-sensitivity}
\end{figure}

All six axes (three angles $\times$ two modalities) exhibit a clear minimum at zero perturbation. For the UAV configuration the roll and pitch curves are approximately symmetric V-shapes climbing from $0.033$ at the optimum to $0.043$--$0.060$ at the $\pm 45^\circ$ extremes; the yaw curve is shallower, reflecting the partial yaw-invariance of a centrally-aimed camera. For the UGV configuration the curves are asymmetric: positive roll and positive pitch produce a much steeper cost rise than the corresponding negative perturbations, because positive perturbations rotate camera $1$ away from the floor slab whereas negative perturbations rotate it further into the slab where redundant cameras already provide coverage. The asymmetry is a direct consequence of the spatially restricted target set under the UGV modality and is exactly the geometric signature one expects for a target-aware optimisation: the gradient of the cost is steep in directions that move coverage off the target set and shallow in directions that move it deeper into the covered region. The absence of any lower-cost configuration within $\pm 10^\circ$ of the optimum, in any axis, confirms that the GA has not merely abandoned the search at an arbitrary point but has converged onto a true local minimum of the cost surface.

\subsection{Optimal Placement Configurations}
\label{sec:sim-optconfig}

This subsection presents the spatial structure of the GA-optimised configurations and contrasts them with the OptiTrack ad-hoc layout. The objective is to expose what the algorithm has learned about the geometry of good placement, and to demonstrate that the choice of target-grid modality (UAV full-volume versus UGV ground-slab) materially shapes the returned layout, as anticipated by Phase~4 of the experimental verification framework (Section~\ref{sec:verification}).

Figure~\ref{fig:configfov} shows the optimised camera poses and field-of-view frustums for both modalities, alongside the OptiTrack ad-hoc rig evaluated on the same target grid. The UAV-optimised layout (Figure~\ref{fig:configfov-uav}, left) distributes the seven cameras across all five mountable faces of the workspace at varying heights, with the three wide-angle cameras seated near the volume corners and the four narrow-angle cameras spread along the long walls. Each camera's frustum interlocks with at least two others over the central capture region. The OptiTrack ad-hoc rig (Figure~\ref{fig:configfov-uav}, right) clusters its seven cameras into two opposing wall groups, leaving the high corners of the volume covered only by single-camera frustums and producing the corner-localised high-cost regions visible in Figure~\ref{fig:costfield-uav}. The combined cost difference, $0.033$ for the GA layout versus $0.269$ for the ad-hoc layout, is consistent with the structural difference between the two: the GA has distributed coverage across the full volume; the ad-hoc rig prioritises the central plane and concedes the periphery.

\begin{figure}[htbp]
    \centering
    \begin{subfigure}[t]{0.49\textwidth}
        \centering
        \includegraphics[width=\textwidth]{Figures/GA/Results/ConfigFOV_GAvsOpti_UAV.pdf}
        \caption{UAV-optimised vs OptiTrack ad-hoc.}
        \label{fig:configfov-uav}
    \end{subfigure}
    \hfill
    \begin{subfigure}[t]{0.49\textwidth}
        \centering
        \includegraphics[width=\textwidth]{Figures/GA/Results/ConfigFOV_GAvsOpti_UGV.pdf}
        \caption{UGV-optimised vs OptiTrack ad-hoc.}
        \label{fig:configfov-ugv}
    \end{subfigure}
    \caption[Camera poses and FOV frustums: GA-optimised vs OptiTrack ad-hoc.]{Side-by-side comparison of the seven-camera GA-best configuration with the OptiTrack ad-hoc baseline for the UAV full-volume target (top) and the UGV ground-slab target (bottom). The GA layout interlocks frustums across the workspace; the ad-hoc layout clusters cameras and concedes coverage at the volume periphery.}
    \label{fig:configfov}
\end{figure}

The UGV-optimised layout (Figure~\ref{fig:configfov-ugv}, left) exhibits a structurally different placement. The cameras are seated lower on the side walls and angled downward, with several positioned close to the ceiling on the long walls but pitched toward the floor; no camera is placed high in the volume aimed horizontally. This concentration of viewing pyramids onto the floor slab is the geometric realisation of the UGV target grid's bias toward $z \leq 0.5$~m: the optimiser has discovered, without explicit instruction, that ceiling-mounted cameras aimed at the floor are the most informative observers for a ground-bound target. The same modality applied to the OptiTrack rig (Figure~\ref{fig:configfov-ugv}, right) yields a combined cost of $0.145$, against $0.018$ for the UGV-optimised layout, an $8\times$ improvement at the same camera count.

The corresponding per-target coverage heatmaps in Figure~\ref{fig:coverage-heatmaps} confirm the spatial implications. The GA-best UAV layout achieves average coverage of $3.51$ cameras per target with $0\%$ of points unobserved and $\sim 68\%$ observed by two or more, against $3.16$, $8.4\%$, and $\sim 73\%$ respectively for the ad-hoc layout (note that the ad-hoc layout's higher fraction of $2{+}$-camera points masks the fact that $8.4\%$ are entirely unobserved). Under the UGV modality the GA layout pushes the unobserved fraction to $0\%$ and the average coverage to $4.13$, against $1.6\%$ unobserved and $4.18$ average for the ad-hoc rig. The headline observation is that the GA can produce a layout that is both more uniform (no zero-coverage holes) and, for the UGV target, more redundantly covered, at no cost to camera count.

\begin{figure}[htbp]
    \centering
    \begin{subfigure}[t]{0.49\textwidth}
        \centering
        \includegraphics[width=\textwidth]{Figures/GA/Results/CoverageHeatmap_GAvsOpti_UAV.pdf}
        \caption{UAV full-volume.}
        \label{fig:coverage-uav}
    \end{subfigure}
    \hfill
    \begin{subfigure}[t]{0.49\textwidth}
        \centering
        \includegraphics[width=\textwidth]{Figures/GA/Results/CoverageHeatmap_GAvsOpti_UGV.pdf}
        \caption{UGV ground slab.}
        \label{fig:coverage-ugv}
    \end{subfigure}
    \caption[Coverage heatmaps for the GA-best and OptiTrack ad-hoc configurations.]{Top row: $XY$ slices of the average number of cameras observing each target point at the slice's $z$. Bottom row: worst-case $XY$ projection (minimum number of cameras observing any target at that $(x,y)$ taken over all $z$). The GA-best layout produces uniform-or-better coverage with no zero-cam holes; the ad-hoc layout exhibits localised under-coverage near the periphery.}
    \label{fig:coverage-heatmaps}
\end{figure}

\paragraph{Modality cross-comparison.}
A salient question for the verification framework is whether a configuration optimised for one modality remains competitive when deployed under the other. Evaluating the UAV-best chromosome on the UGV grid yields a combined cost of $\approx 0.034$ (against $0.018$ for the UGV-best chromosome at the same grid), and the UGV-best chromosome evaluated on the UAV grid yields $\approx 0.072$ (against $0.033$ for the UAV-best). Both cross-evaluations remain well below the OptiTrack ad-hoc cost on either grid, but each is roughly $2\times$ above the modality-matched best. This confirms the design hypothesis of Section~\ref{sec:cost-function}: a sufficient placement for one operating context is not optimal for the other, and the discretisation modality is a meaningful encoding of deployment intent rather than an arbitrary internal parameter. The physical verification (Phase~4 of Section~\ref{sec:verification}) is therefore expected to find that the modality-matched configuration outperforms the cross-matched configuration on the corresponding tracking task, with the cross-matched layout still outperforming the ad-hoc baseline.

\paragraph{Summary of best returned costs.}
Table~\ref{tab:best-costs} aggregates the best returned cost across all $80$ runs per cell, alongside the ad-hoc baseline at the same conditions, so that the verification chapter can refer to the headline numbers in one place.

\begin{table}[htbp]
\centering
\caption[Headline best costs across all cells at $k=7$, $\Delta=1$~m.]{Headline best returned combined cost $J$ across all runs in each $(k=7,\Delta=1~\mathrm{m})$ cell, together with the OptiTrack ad-hoc baseline evaluated at the same conditions and the corresponding improvement factor.}
\label{tab:best-costs}
\small
\begin{tabular}{l l c c c}
\toprule
TT & GM & GA best $J$ & OptiTrack ad-hoc $J$ & Improvement factor \\
\midrule
UAV & Uniform & 0.0330 & 0.2691 & $8.2\times$ \\
UAV & Normal  & 0.0281 & 0.2393 & $8.5\times$ \\
UGV & Uniform & 0.0179 & 0.1449 & $8.1\times$ \\
UGV & Normal  & 0.0126 & 0.1218 & $9.7\times$ \\
\bottomrule
\end{tabular}
\end{table}

Across all four cells the GA achieves an order-of-magnitude reduction in the composite cost relative to the ad-hoc OptiTrack arrangement, with the improvement most pronounced under the normal-concentrated UGV modality, where the optimiser most aggressively concentrates coverage on the floor slab. Whether these simulation-level reductions translate into proportionate gains in measured calibration residual, marker-reconstruction completeness, and static accuracy in the physical OptiTrack deployment is the question that the experimental protocol of Section~\ref{sec:verification} is designed to answer.
